\section{Divergence of the prime reciprocals}\label{sec:divergence}

% @claim C001 from=thm:divergence
\begin{thm}\label{thm:divergence}
The series $\sum_{p \in \primes} 1/p$ diverges.
\end{thm}

The proof compares two estimates of $\log \zeta(s)$ as $s$ decreases to
one. The first comes from the integers, the second from the primes.

\begin{lem}\label{lem:log}
Let $s > 1$. Then:
\begin{enumerate}
  \item\label{it:lower} $\log \zeta(s) \ge \log \frac{1}{s-1}$;
  \item\label{it:upper} $\log \zeta(s) \le \sum_{p} p^{-s} + C$, where
  $C = \sum_p p^{-2s} \le \zeta(2)$.
\end{enumerate}
\end{lem}

\begin{proof}
\textbf{Step 1.} For the lower bound compare the sum with an integral:
\begin{align}
  \zeta(s) &\ge \int_1^\infty t^{-s}\,dt \label{eq:integral}\\
           &= \frac{1}{s-1}. \label{eq:integral-value}
\end{align}

\textbf{Step 2.} For the upper bound take logarithms in
\cref{eq:euler-product} and use $-\log(1-x) \le x + x^2$ on
$0 \le x \le 1/2$, which holds since

\[
  -\log(1-x) = \sum_{k \ge 1} \frac{x^k}{k}
    \le x + \frac{x^2}{2} \sum_{k \ge 0} x^k.
\]

and the last sum is at most two.
\end{proof}

\begin{proof}[Proof of \cref{thm:divergence}]
Suppose the series converged, to $S$ say. By \cref{lem:log}~\ref{it:upper},
$\log \zeta(s) \le S + \zeta(2)$ for every $s > 1$.

By
\cref{lem:log}~\ref{it:lower} the left side tends to infinity as
$s \to 1^+$
\[
  \log \frac{1}{s-1} \to \infty,
\]
a contradiction.
\end{proof}

\begin{cor}\label{cor:infinite}
There are infinitely many primes.
\end{cor}

\begin{proof}
A finite set of primes has a finite sum of reciprocals; apply
\cref{thm:divergence}.
\end{proof}

The two bounds \crefrange{eq:integral}{eq:integral-value} are the whole of
the lower estimate, and \eqref{eq:integral-value} is sharp as $s \to 1^+$.

\subsection{A sharper estimate}\label{sec:sharper}

A sharper statement holds: the partial sums grow like $\log\log x$.

\begin{claim}\label{claim:mertens}
$\sum_{p \le x} 1/p = \log\log x + O(1)$.
\end{claim}

We do not prove \cref{claim:mertens}; Mertens' proof~\cite{mertens1874}
uses the same comparison with more care, and \enquote{the constant that
appears is not a simple one}~\cite{primer2001}. The steps are:
\begin{enumerate}
  \item \emph{(Chebyshev's bound.)} The primes up to $x$ number $O(x/\log x)$.
  \item \emph{(Partial summation.)} The sum over primes becomes an integral
  against the counting function.
\end{enumerate}

\begin{tcolorbox}
A computation of the partial sums for $x \le 10^6$ agrees with
$\log\log x + 0.2615$ to three decimals.
\end{tcolorbox}

The computation's record has the form
\begin{verbatim}
% @claim C999 from=claim:mertens
x = 1000000, sum = 2.887328
\end{verbatim}

\section{Remarks}\label{sec:remarks}

\begin{rem}
Euler's original argument multiplies divergent series as if they
converged; the inequalities of \cref{lem:log} are what make it rigorous,
and \thmref{divergence} is what it was after.
\end{rem}

\begin{proof}[Proof of the claims in \cref{sec:product}]
Each claim there follows from absolute convergence.
\end{proof}

\begin{lem}\label{lem:tail}
For $s \ge 2$, $\sum_{n > x} n^{-s} \le 1/x$.
\end{lem}
The bound is crude.
\begin{proof}
Compare with $\int_x^\infty t^{-2}\,dt$.

The integral is $1/x$.
\end{proof}

\ifdefined\eulerconst We write $\eulerconst$ for Euler's constant.\fi

\subsubsection*{Further reading}
\Cref{sec:sharper} and the section \nameref{sec:remarks} point to
\citet{primer2001}; \vref{lem:tail} is used nowhere else.

\section*{Acknowledgements}
We thank the reader for reading.
